\documentclass[14pt]{extreport}
\usepackage{indentfirst}
\usepackage[T2A]{fontenc}
\usepackage[utf8]{inputenc}
\usepackage[russian, english]{babel}
%\setmainfont{Times New Roman}
%\setsansfont{Arial}
%\setmonofont{Courier New}
\addto\captionsrussian{%
	\renewcommand{\contentsname}{Содержание}%
}
\addto\captionsrussian{%
	\renewcommand{\bibname}{Список литературы}%
}
\usepackage[top = 20mm, left = 30mm, right = 15mm, bottom = 20mm]{geometry}
\setlength\parindent{12.5mm}
\renewcommand{\baselinestretch}{1.5}
\def\hrf#1{\hbox to#1{\hrulefill}}
\usepackage{ragged2e}
\usepackage{caption}
\usepackage{multicol}
\usepackage{mathtools}
\usepackage{url}
\usepackage{xurl}
\usepackage{amsmath,amsfonts,amssymb}
\usepackage{courier}
\usepackage{hyperref}
\usepackage{enumitem}
\usepackage{listings}
\usepackage{xcolor}
\usepackage{courier}
\usepackage{verbatim}
\usepackage{graphicx}
\usepackage{float}

\renewcommand\thesection{\arabic{section}}

% Цвета для подсветки
\definecolor{bluekeywords}{rgb}{0.125,0.375,0.625}
\definecolor{greencomments}{rgb}{0.125,0.5,0.125}
\definecolor{redstrings}{rgb}{0.875,0.125,0.125}
\definecolor{graynumbers}{rgb}{0.5,0.5,0.5}
\definecolor{lightgray}{gray}{0.95}
\definecolor{codegreen}{rgb}{0,0.6,0}
\definecolor{codegray}{rgb}{0.5,0.5,0.5}
\definecolor{codepurple}{rgb}{0.58,0,0.82}
\definecolor{backcolour}{rgb}{0.95,0.95,0.95}

% Макрос для кода (Courier + серый фон)
%\usepackage{tgcursor}  % делает \texttt = TeX Gyre Cursor (аналог Courier)
\newcommand{\code}[1]{\colorbox{lightgray}{\texttt{#1}}}


\lstset{
	language=C++,
	basicstyle=\ttfamily\small,
	keywordstyle=\color{bluekeywords}\bfseries,
	commentstyle=\color{greencomments},
	stringstyle=\color{redstrings},
	numberstyle=\color{graynumbers},
	numbers=left,
	numbersep=10pt,
	columns=fullflexible,
	tabsize=4,
	showspaces=false,
	showstringspaces=false,
	breaklines=true,
	backgroundcolor=\color{lightgray},
	frame=single,
	rulecolor=\color{gray!30!black},
	captionpos=b
}


\begin{document}
	\sloppy
	\thispagestyle{empty}
	\centering{
		Министерство образования и науки Российской Федерации\\
		Федеральное государственное автономное образовательное учреждение высшего образования\\
		«Санкт-Петербургский политехнический университет Петра Великого»\\[10pt]
		Институт компьютерных наук и кибербезопасности\\
		Высшая школа технологий искусственного интеллекта\\
		02.03.01 Математика и компьютерные науки\\
		\vspace{\fill}
		{\large
			Отчёт по дисциплине\\
			{\Large \bf Вычислительная математика\\Лабораторная работа №3}\\
			{\Large <<Кубические сплайны>>}
            \\Вариант №15
			\vspace{\fill}
		}
		
		\flushright
		{
			{\bf Выполнил:}\\
			студент группы 5130201/40003\\
			Гумницкий К. С.\\
			{\bf Принял:}\\
			В. Г. Пак\\
			<<\hrf{2em}>> \hrf{6em} 2026~г.
		}
		
		\centering{
			Санкт-Петербург, 2026
		}

        \newpage

        \justifying

        \section{Задание}
        Построить для данной табличной функции интерполяционный кубический сплайн дефекта 1 с заданными граничными условиями. Коэффициенты вычислять с четырьмя знаками после запятой в мантиссе. Проверить все условия, налагаемые на сплайн. 

        \section{Цель работы}
        Научиться строить интерполяционные кубические сплайны для табличных функций

        \section{Ход работы}
        \begin{table}[h!]
\centering
\begin{tabular}{|c|c|}
\hline
\multicolumn{2}{|c|}{\textbf{Вариант 15}} \\ \hline
\multicolumn{2}{|c|}{Граничные условия: отсутствие узла} \\ \hline
$x_i$ & $y_i$ \\ \hline
1,000 & 0,0000 \\ \hline
1,510 & 0,2729 \\ \hline
2,100 & 0,3533 \\ \hline
2,750 & 0,3679 \\ \hline
3,420 & 0,3595 \\ \hline
3,915 & 0,3486 \\ \hline
4,350 & 0,3380 \\ \hline
4,800 & 0,3268 \\ \hline
5,200 & 0,3171 \\ \hline
\end{tabular}
\end{table}
\code{}

Функция \code{spline} реализует построение кубического интерполяционного сплайна дефекта 1 с граничными условиями «отсутствие узла». Пусть задана сетка узлов $x_0 < x_1 < \ldots < x_n$ и соответствующие им значения функции $y_i = f(x_i)$. На каждом частичном интервале $[x_i, x_{i+1}]$ сплайн представляется кубическим полиномом $S_i(x) = a_i + b_i(x-x_i) + c_i(x-x_i)^2 + d_i(x-x_i)^3$, где коэффициенты $a_i$, $b_i$, $c_i$, $d_i$ подлежат определению. Из условия интерполяции непосредственно получаем $a_i = y_i$. Для нахождения остальных коэффициентов используются условия непрерывности сплайна, его первой и второй производных во внутренних узлах, что приводит к системе линейных уравнений относительно вторых производных $c_i$ (так называемых моментов). Для $i = 1, 2, \ldots, n-1$ выполняются соотношения

\[
h_{i-1} \cdot c_{i-1} + 2(h_{i-1} + h_i) \cdot c_i + h_i \cdot c_{i+1} = 3\left(\frac{y_{i+1}-y_i}{h_i} - \frac{y_i-y_{i-1}}{h_{i-1}}\right),
\]

где $h_i = x_{i+1} - x_i$ — шаг сетки. Граничные условия «отсутствие узла» означают, что третья производная сплайна непрерывна в точках $x_1$ и $x_{n-1}$, что эквивалентно равенствам $c_1 = c_0$ и $c_n = c_{n-1}$. Для левого края это приводит к уравнению

\[
-h_1 \cdot c_0 + (h_0 + h_1) \cdot c_1 - h_0 \cdot c_2 = 0,
\]

а для правого края — к уравнению

\[
-h_{n-1} \cdot c_{n-2} + (h_{n-2} + h_{n-1}) \cdot c_{n-1} - h_{n-2} \cdot c_n = 0.
\]

Таким образом формируется замкнутая система $(n+1)$ линейных уравнений относительно $(n+1)$ неизвестных $c_0, c_1, \ldots, c_n$, которая решается методом Гаусса или любой другой подходящей численной схемой. После нахождения всех $c_i$ коэффициенты $b_i$ и $d_i$ вычисляются по явным формулам:

\[
b_i = \frac{y_{i+1}-y_i}{h_i} - \frac{h_i(2c_i + c_{i+1})}{3}, \qquad
d_i = \frac{c_{i+1} - c_i}{3h_i}.
\]

Ниже приведён код функции сплайн

    \begin{lstlisting}[language=Python]
def spline(x, y):
    x = np.array(x, dtype=float)
    y = np.array(y, dtype=float)
    n = len(x) - 1
    
    h = np.diff(x)
    
    a = y[:-1].copy()
    
    A = np.zeros((n + 1, n + 1))
    rhs = np.zeros(n + 1)
    
    for i in range(1, n):
        A[i, i-1] = h[i-1]
        A[i, i] = 2 * (h[i-1] + h[i])
        A[i, i+1] = h[i]
        rhs[i] = 3 * ((y[i+1] - y[i]) / h[i] - (y[i] - y[i-1]) / h[i-1])
    
    if n >= 2:
        A[0, 0] = -h[1]
        A[0, 1] = h[0] + h[1]
        A[0, 2] = -h[0]
        rhs[0] = 0
        
        A[n, n-2] = -h[n-1]
        A[n, n-1] = h[n-2] + h[n-1]
        A[n, n] = -h[n-2]
        rhs[n] = 0
    else:
        A[0, 0] = 1
        rhs[0] = 0
        A[n, n] = 1
        rhs[n] = 0
    
    c = np.linalg.solve(A, rhs)
    
    b = np.zeros(n)
    d = np.zeros(n)
    
    for i in range(n):
        b[i] = (y[i+1] - y[i]) / h[i] - h[i] * (2 * c[i] + c[i+1]) / 3
        d[i] = (c[i+1] - c[i]) / (3 * h[i])
    
    return a, b, c, d
		\end{lstlisting}

        \newpage
        
        \section{Результаты}
        

\newpage

        \section{Приложение}
\lstinputlisting[language=Python,caption={Файл Lab3.py}]{Lab3.py}

\end{document}